\section{ПРОЕКТИРОВАНИЕ И РАЗРАБОТКА МЕДИА-РЕСУРСА}

\subsection{Обоснование архитектуры и выбора технологического стека}
Для реализации информационной платформы проекта «ЭкоМедиа» требовалось создать надежный, быстрый и легко доступный для целевой аудитории веб-сайт. При выборе технологического стека учитывались такие критерии, как скорость загрузки страниц, кроссбраузерность, простота развертывания (хостинга) и независимость от сложных серверных инфраструктур. 

В результате анализа было принято решение отказаться от тяжеловесных систем управления контентом (CMS) и серверных фреймворков в пользу концепции статического сайта (Static Site Generation / Plain HTML). Архитектура статического сайта подразумевает, что веб-страницы генерируются и сохраняются в виде готовых HTML-файлов, которые отдаются пользователю веб-сервером без необходимости обращения к базе данных или выполнения серверных скриптов при каждом запросе.

Преимущества выбранного подхода:
\begin{enumerate}
    \item \textbf{Высокая производительность.} Отсутствие времени на серверную генерацию страниц обеспечивает мгновенную загрузку ресурса, что критически важно для удержания внимания молодежной аудитории.
    \item \textbf{Безопасность.} Исключение серверной логики и баз данных делает сайт практически неуязвимым для классических веб-атак, таких как SQL-инъекции или межсайтовый скриптинг (XSS).
    \item \textbf{Простота хостинга.} Статические файлы легко размещаются на бесплатных платформах, таких как GitHub Pages, что идеально подходит для учебного проекта и обеспечивает надежную работу 24/7.
\end{enumerate}

Основным технологическим стеком для фронтенд-разработки выступили стандарты HTML5 и CSS3. Для управления версиями кода и совместной работы была использована система Git и удаленный репозиторий GitHub.

\subsection{Разработка пользовательского интерфейса и дизайна}
Дизайн веб-сайта играет ключевую роль в восприятии экологического контента. Традиционно ресурсы по экологии оформляются в утилитарном или чрезмерно строгом стиле. Для проекта «ЭкоМедиа» была выбрана современная концепция \textbf{Glassmorphism} (эффект матового стекла) в сочетании с темной темой оформления (Dark Mode).

Темная тема была выбрана не случайно: она тесно перекликается с темой цифровой гигиены. Использование темного фона (оттенок Slate — \#0f172a) снижает излучение света экраном, что уменьшает утомляемость глаз и сокращает подавление выработки мелатонина при просмотре сайта в вечернее время. В качестве акцентного цвета был выбран изумрудный зеленый (Emerald — \#10b981), символизирующий жизнь, природу и экологию, а также небесно-голубой цвет воды.

Эффект Glassmorphism достигается за счет применения свойств CSS \texttt{backdrop-filter: blur(10px)} и полупрозрачных фонов с альфа-каналом (RGBA). Это позволяет создать ощущение глубины и многослойности интерфейса, не перегружая его лишними графическими элементами. Карточки контента выглядят как стеклянные пластины, парящие над плавно переливающимся градиентным фоном.

Особое внимание при проектировании уделялось адаптивности (Responsive Web Design). Использование медиазапросов (\texttt{@media}) и технологии CSS Flexbox/Grid позволило обеспечить корректное и удобное отображение контента на любых устройствах: от широкоформатных мониторов до мобильных телефонов.

\subsection{Структура сайта и навигация}
Информационная архитектура сайта спроектирована с учетом максимального удобства пользователя. Навигационная панель (header) закреплена в верхней части экрана (sticky positioning) и всегда остается доступной при прокрутке страницы.

Сайт состоит из следующих основных разделов:
\begin{itemize}
    \item \textbf{Главная страница (index.html):} Выполняет функцию лендинга, содержит крупные акцентные заголовки (hero section) и призыв к действию, мотивирующий погрузиться в изучение проекта.
    \item \textbf{О проекте (about.html):} Подробное описание миссии, актуальности проблемы, целей и задач исследования в области экологии человека.
    \item \textbf{Участники (participants.html):} Презентация команды проекта. Раздел содержит карточки участников с описанием их конкретного вклада в реализацию (менеджмент, дизайн, создание сценариев, SMM).
    \item \textbf{Журнал (journal.html):} Формат блога или таймлайна, где задокументированы ключевые этапы развития проекта: от первоначального планирования и создания контента до анализа откликов аудитории.
    \item \textbf{Ресурсы (resources.html):} База знаний с полезными ссылками на внешние авторитетные источники, такие как доклады ВОЗ и профильные экологические медиа.
    \item \textbf{Telegram-бот (bot.html):} Специальная посадочная страница (Landing Page) для продвижения разработанного программного продукта, содержащая описание функционала и прямую ссылку на мессенджер.
\end{itemize}

Весь контент оптимизирован для быстрого восприятия (сканирования): тексты разбиты на небольшие абзацы, активно используются маркированные списки, выделения жирным шрифтом и визуальные иконки.
